\homework{3}{Tue 06 Sep 2022 14:42}{Week 3 due Sep 14}

Herstein, section 2.4, page 64 : Questions 8,13,16,18,27,31,37 
\begin{enumerate}
	\item 8. If every right coset of $H$ in $G$ is a left coset of $H$ in $G$, prove that $a H a^{-1}=H$ for all $a \in G$
\begin{proof}
	We interpret the assumption as: For any right coset $Ha$ in $G$,  $Ha = aH$. Note that this also implies every left coset of $H$ in $G$ is a right coset. Now

	\begin{align*}
		x &\in aHa^{-1} \\
		  &\iff x \in a (Ha^{-1} )\\
		  &\iff x \in a (a^{-1} H)\\
		  &\iff x \in (aa^{-1} )H \\
		  &\iff x \in H
	.\end{align*}
\end{proof}

	\item 13. Find the orders of all the elements of $U_{18}$. Is $U_{18}$ cyclic?
\begin{proof}
	$U_{18} = \{[1],[5],[7],[11],[13],[17]\} $.
We find order of each element:
\begin{align*}
	\left\langle 1 \right\rangle &= \{[1]\}  \\
	\left\langle [5] \right\rangle &= \{[5],[7],[17],[13],[11],[1]\}  \\
	\left\langle [7] \right\rangle &= \{[7],[13],[1]\}  \\
	\left\langle 11 \right\rangle &= \{[11],[13],[17],[7],[5],[1]\}  \\
	\left\langle 13 \right\rangle &= \{[13],[7],[1]\}  \\
	\left\langle 17 \right\rangle &= \{[17],[1]\}
.\end{align*}

Hence, they have order $1, 6, 3, 6, 3, 2$ respectively. The group $U_{18}$ is hence cyclic since it is generated by $[5]$ and $[11]$.
\end{proof}

	\item 16. If $G$ is a finite abelian group and $a_1, \ldots, a_n$ are all its elements, show that $x=a_1 a_2 \cdots a_n$ must satisfy $x^2=e$.
\begin{proof}
	Without loss of generality we assume $a_1, \ldots, a_n$ are distinct.
	Let $a_1'a_2'\ldots a_n'$ be a permutation of $a_1a_2\ldots a_n$. Because $G$ is finite and abelian, permutation does not change the product. Hence $x^2 = (a_1a_2\ldots a_n)(a_1a_2 \ldots a_n) = (a_1a_2\ldots a_n)(a_1'a_2' \ldots a_n')$.

	Now by the inverse axiom, there must exist some $a \in G$ such that $a_1 a = e$. Without loss of generality we can call it $a_1'$. Hence
\[
	(a_1a_2 \ldots a_n) (a_1'a_2' \ldots a_n') = 
	(a_1a_1')(a_2\ldots a_n) (a_2' \ldots a_n') = 
	(a_2 \ldots a_n) (a_2' \ldots a_n')
.\] 
Then we consider the inverse of $a_2$. It cannot be $a_1'$, otherwise $a_2 = \left( a_1' \right) ^{-1} = \left( a_1^{-1} \right) ^{-1} = a_1$, a contradiction. Without loss of generality we may call its inverse $a_2' = a_2^{-1}$. Now
\[
	(a_2 \ldots a_n) (a_2' \ldots a_n') = 
	(a_2 a_2') (a_3 \ldots a_n) (a_3' \ldots a_n')= 
	(a_3 \ldots a_n) (a_3' \ldots a_n')
.\] 
Continue in this fashion we can cancel all terms and onclude that $x^2 = e$.
\end{proof}
\begin{proof}
	(simpler)

	The inverse  \begin{align*}
		i: G &\longrightarrow G \\
		a &\longmapsto i(a) = a^{-1}
	\end{align*} is a bijection. 
	Since $G$ is abelian,
	\[
		(a_1a_2\ldots a_n) =
		(a_1^{-1} a_2^{-1} \ldots a_n^{-1})
	.\] 

	Now
	\begin{align*}
		x^2 &= (a_1a_2\ldots a_n)(a_1a_2\ldots a_n)\\
		&= (a_1a_2\ldots a_n) (a_1^{-1} a_2^{-1} \ldots a_n^{-1})\\
		&= (a_1a_1^{-1})\ldots(a_n a_n^{-1}) \\
		&= e
	.\end{align*} 

\end{proof}

	\item 18. Using the results of Problems 15 and 16 , prove that if $p$ is an odd prime number, then $(p-1) ! \equiv-1 \bmod p$. (This is known as Wilson's Theorem.) It is, of course, also true if $p=2$.
\begin{proof}
	Consider the group $U_p$. Since $p$ is prime, the elements of the group are exactly the natural numbers from $1$ up to $p -1$. Now
	\[
		(p-1)! = 1 \cdot 2 \cdot \ldots \cdot (p-1) =  a_1 a_2 \ldots a_n = x
	.\] 
By Exercise 16, $x^2 = 1$. By Exercise 15, any member of  $U_p$ that is not $1$ or $p -1$ does not satisfy $x^2 \equiv_p 1$. Note that: \[
	x^2 \equiv_p 1 \iff x^2 = e \iff x = x^{-1} 
.\] 
Therefore all members except $1 $ or $p -1$ are not their own inverse. Hence they must be pairs that are inverse of each other. Because $U_p$ has even order, and there are even number of elements excluding $1$ and $p -1$, the remaining elements can all pair up correctly. Because  $U_p$ is abelian, we may group the pairs together and cancel them out. Now
\[
	(p-1)! = 1 \cdot (p-1) \equiv_p -1
.\] 
	
\end{proof}

	\item 27. If $a H=b H$ forces $H a=H b$ in $G$, show that $a H a^{-1}=H$ for every $a \in G$.
\begin{proof}
	Take $a \in G$, $b \in H$. We want to show $ah \in Ha$. 
	\begin{claim}
		There exists $b \in G$ such that $a h \in b H$ and $b H = H b$.
	\end{claim}
\end{proof}

	\item 31. If $o(a)=m$ and $a^s=e$, prove that $m \mid s$.
	\item 37. In a cyclic group of order $n$, show that for each positive integer $m$ that divides $n$ (including $m=1$ and $m=n$ ) there are $\varphi(m)$ elements of order $m$.
\begin{proof}
	An element $a^k \in \left\langle a \right\rangle = G$ has order $m$ iff
	\begin{enumerate}
		\item $(a^k)^m = a^{km} = e$. In other words, $n | k m$.
		\item $m$ is the smallest positive integer such that $ km $ is a multiple of $n$. In other words, for any $p < m$, we have $n \nmid kp$.
	\end{enumerate}
	Note that $1\le m\le n$ (since $m$ divides $n $), and $0\le k\le n-1$ (since $a^k$ are distinct elements of $G$).

	We introduce a claim that is central to the proof , which we prove later:
	\begin{claim}
		The two conditions above are met, if and only if \[
			(k,n) = \frac{n}{m}
		.\] 
	\end{claim}

	With the claim proved, since
	\[
		\left( k, \frac{n}{n / m} \right) = \left( k,m \right) = 1 
	,\] 
	we have $o(a^k) = m \iff (k,m) = 1$, and there are $\varphi(m)$ choices of $k$ such that $(k,m) = 1$.
\end{proof}
	\begin{notation}
		For integer $m$ and prime $x$, write  $d_x(m)$ for the \textit{multiplicity} of $x$ in the unique prime factorization of $m$.
	\end{notation}
	\begin{lemma}
		(Properties of Multiplicity) The following statements hold:
		\begin{enumerate}
			\item $d_x(nm) = d_x(n) + d_x(m)$ 
			\item $d_x(\frac{n}{m}) = d_x(n) - d_x(m)$ if $m | n$.
			\item  $d_x((m,n)) = \min(d_x(m), d_x(n))$.
			\item If $m | n$, then one has $d_x(m) \le  d_x(n)$. Conversely, if $d_x(m) > d_x(n)$ then $m \nmid n$.
		\end{enumerate}
	\end{lemma}
	\begin{proof}
		
	\end{proof}
	
	\begin{lemma}
	If $p < m$, then there exists a prime $x>1$ such that $x|m $ and $d_x(m) > d_x(p)$.
	\end{lemma}
	\begin{proof}
		
	\end{proof}
	
	
	\begin{proof}[Proof of claim]
	Suppose $(k,n) = \frac{n}{m}$. Hence $(km, nm) = n$. This means that $n |km $, satisfying condition (a).  

	Suppose there exists positive integer $p$ such that  $ p<m$ and $n | k p$. Use the above lemma to choose prime $x$ such that  $x>1$, $x|m$ and $d_x(m) > d_x(p)$. We claim that \[
		d_x(n) = d_x(m) + d_x(k)
	.\] 
	To verify this, since $(k,n) = \frac{n}{m}$, we have \[
		\min\left( d_x(k), d_x(n) \right) = d_x(n) - d_x(m)
	.\] 
	Do some arithmetic manipulations and obtain \[
		d_x(m) = \max(0,d_x(n) - d_x(k))
	.\] 
	Since $x | m$, $d_x(m) \ge 1$, therefore \[
		d_x(m) = d_x(n) - d_x(k)
	.\] 

	Now \[
		d_x(pk) = d_x(p) + d_x(k) < d_x(m) + d_x(k) = d_x(n)
	.\] 
	Hence $n \nmid pk $.

	For the other direction, assume that $a^k$ has order $m$ and hence satisfies the two conditions. We first introduce
	\begin{lemma}
		For any positive integer $n,k,m$, we have \[
			n|km \implies n|(k,n) m
		.\] 
	\end{lemma}
	\begin{proof}
		TBD
	\end{proof}
	

	Then since $n | km$, we apply the previous lemma and conclude that $n | (k,n) m$. Hence there exists positive integer $c$ such that $cn = (k,n) m$. We claim that $c = 1$, which will conclude the proof.

	To verify that $c=1$, suppose for the sake of contradiction that $c > 1$. Now $(k,n) = c \frac{n}{m}$ since $m | n$. Hence  $c$ is a factor of both $k$ and $n$, and we have \[
		\left( \frac{k}{c}, \frac{n}{c} \right) = \frac{n}{m}
	.\] 
	Hence in particular \[
		\frac{n}{m} \mid \frac{k}{c}
	.\] 
	Hence \[
		n \mid \frac{k}{c} m
	\]
	where $\frac{k}{c}<k$, a contradiction to condition 2.
\end{proof}


\end{enumerate}
Herstein, section 2.5, page 73 : Questions 2,6,7
